\section{L’atome d’hydrogène}

    \subsection{Modèle de Bohr}

        Ce modèle, bien que faux, utilise la théorie classique pour donner les valeurs exactes des niveaux d’énergie de l’atome d’hydrogène, et possède une importance historique notable. 

        \textbf{Modèle :} L’électron a une trajectoire circulaire autour d’un proton fixe qui sert d’origine, on considère des interactions coulombiennes.

        Dans ce cadre, d’après le PFD, 
        \[ m_e \frac{v^2}{r} = \frac{e^2}{r^2} \quad \text{avec} \quad e^2 = \frac{q^2}{4 \pi \varepsilon_0} \]   
        En postulant que la quantification du moment cinétique $L_z = n \hbar$, on doit avoir 
        \begin{align*}
            &L_z = r m_e v = n \hbar  \\
            \implies &v_n = \frac{e^2}{n\hbar} = \alpha \frac{c}{n} \\
            \implies &r_n = \frac{n^2}{a_1} \\
            \implies &E_n = \frac{1}{2} m_e v_n^2 - \frac{e^2}{r_n} = - \frac{E_I}{n^2} \\
        \end{align*}
        Les constantes apparaissant ici sont :
        \begin{itemize}
            \item $\alpha= \frac{e^2}{\hbar c } \approx \frac{1}{137}$, la \textbf{constante de structure fine}.
            \item $a_1 = \frac{\hbar^2}{m_e e^2} \approx 53\text{pm}$, le \textbf{rayon de Bohr}.
            \item $E_I = \frac{m_e e^4}{2 \hbar^2} \approx 13.6\text{ev}$, l’\textbf{énergie d’ionisation} de l’atome d’hydrogène.
        \end{itemize}
        Ce résultat établit par Bohr explicite bien la quantification des niveaux d’énergie, dont les fréquences de transition sont 
        \[ h \nu = E_{n_2} - E_{n_1} = E_I \left(\frac{1}{n_1^2} - \frac{1}{n_2^2}\right) \]   
        Ces résultats conformes aux attentes expérimentales sont toutefois loin d’être convenables, étant donné qu’un problème à une échelle si petite doit absolument être traité de façon quantique.

    \subsection{Mouvement dans un potentiel central}

        \subsubsection{Cas général}

    Considèrons désormais le problème dans le cadre de la mécanique quantique : deux particules de masses $m_1$ et $m_2$ interagissent entre elles selon un potentiel dépendant uniquement de la distance $r = \nor{\vec{r_2} - \vec{r_1}}$ entre les deux particules. Le système est alors invariant par translation et par rotation. On choisira aussi l’origine des énergies afin que $V(r) \underset{r \to \infty}{\longrightarrow} 0$. L’hamiltonien du système est alors 
    \[ \hat{H}_{\text{tot}} = \frac{\hat{p}_1^2}{2m_1} + \frac{\hat{p}_2^2}{2 m_2} + V(r) \]   

    \begin{omed}{Simplification}
        On peut se ramener à l’étude de deux systèmes indépendants, correspondant d’une part au mouvement libre du centre de masse, de masse $M = m_1 + m_2$, et d’autre part au mouvement d’une particule fictive plongée dans le potentiel $V(r)$ de masse réduite $\mu$.
    \end{omed}

    \begin{demo}{Démonstration}
        On effectue un changement de variable en posant 
        \[ \vec{R} = \frac{m_1 \vec{r_1} + m_2 \vec{r_2}}{M} \]   
        qui correspond à la position du centre de masse. Ce changement de variable nous permet de réinterpréter $\vec{r_1}$ et $\vec{r_2}$ comme
        \begin{align*}
            \vec{r_1} &= \vec{R} - \frac{\mu}{M} \vec{r} \\
            \vec{r_2} &= \vec{R} + \frac{\mu}{M} \vec{r} \\
        \end{align*}
        où $\mu$ est la masse réduite définie par 
        \[ \frac{1}{\mu} = \frac{1}{m_1} + \frac{1}{m_2} \]   
        On peut alors introduire une nouvelle fonction d’onde de la forme 
        \[ \Phi(\vec{R}, \vec{r})= \eta \Psi \left(\vec{R} - \frac{\mu}{M} \vec{r}, \vec{R} + \frac{\mu}{M} \vec{r}\right) \]
        où $\eta$ est une constante de normalisation de la nouvelle fonction d’onde. On se place ainsi dans un nouvel espace de Hilbert $\mathcal{E}_H = \mathcal{L}^2(\mathbf{R}^3) \otimes \mathcal{L}^2(\mathbf{R}^3)$ où le premier espace correspond à la position $\vec{R}$ du centre de masse et le second à la position $\vec{r}$ de la particule fictive de masse $\mu$. Identifions désormais les observables impulsion $\hat{\vec{P}}$ et $\hat{\vec{p}}$ associées à ces deux nouvelles variables, en utilisans la définition de l’impulsion comme générateur infinitésimal du groupe des translations :
        \begin{align*}
            \Phi(\vec{R} - d\vec{a}, \vec{r}) 
            &= \Phi(\vec{R}, \vec{r}) - \eta \frac{\partial \Phi}{\partial \vec{r}_1} \cdot d\vec{a} - \eta \frac{\partial \Phi}{\partial \vec{r}_2} \cdot d\vec{a} \\
            &= \left(1 - \frac{i}{\hbar}(\hat{\vec{p}}_1 + \hat{\vec{p}}_2) \cdot d\vec{a}\right) \Phi(\vec{R}, \vec{r}) \\
        \end{align*}
        On en déduit que 
        \[ \hat{\vec{P}} = \hat{\vec{p}}_1 + \hat{\vec{p}}_2 \] 
        De la même façon, 
        \begin{align*}
            \Phi(\vec{R}, \vec{r} - d\vec{a})
            &= \Phi(\vec{R}, \vec{r}) + \eta \frac{\mu}{m_1} \frac{\partial \Phi}{\partial \vec{r}_1} \cdot d\vec{a} - \eta \frac{\mu}{m_2}\frac{\partial \Phi}{\partial \vec{r}_2} \cdot d\vec{a} \\
            &= \left(1 - \frac{i}{\hbar} \mu \left[\frac{\hat{\vec{p}}_2}{m_2} - \frac{\hat{\vec{p}}_1}{m_1}\right] \cdot d\vec{a}\right) \Phi(\vec{R}, \vec{r}) \\
        \end{align*}
        L’impulsion associée à la variable spaciale $\vec{r}$ est donc
        \begin{align*}
            \hat{\vec{p}} 
            &= \mu \left[\frac{\hat{\vec{p}}_2}{m_2} - \frac{\hat{\vec{p}}_1}{m_1}\right] \\
            &= \frac{m_1 \hat{\vec{p}}_2 - m_2 \hat{\vec{p_1}}}{M} \\
        \end{align*}
        On peut ainsi réexprimer les impulsions des deux particules en fonction des nouvelles observables :
        \begin{align*}
            \hat{\vec{p}}_1 &= \frac{m_1}{M} \hat{\vec{P}} - \hat{\vec{p}} \\
            \hat{\vec{p}}_2 &= \frac{m_2}{M} \hat{\vec{P}} + \hat{\vec{p}} \\
        \end{align*}
        En particulier, cela permet de réexprimer l’énergie cinétique totale du système puisque 
        \[ \frac{\hat{p}_1^2}{2m_1} + \frac{\hat{p}_2^2}{2 m_2} = \frac{\hat{P}^2}{2M} + \frac{\hat{p}}{2 \mu} \]
        L’hamiltonien du système devient alors 
        \[ \hat{H}_{\text{tot}} = \hat{H}_{\text{cm}} + \hat{H} \]   
        où $\hat{H}_{\text{cm}} = \frac{\hat{P}^2}{2M}$ est l’hamiltonien correspondant au centre de masse et $\hat{H} = \frac{\hat{p}^2}{2\mu} + V(\hat{r})$ est l’hamiltonien correspondant à la particule fictive placée dans le potentiel central $V(r)$. Ces hamiltoniens commutent car ils agissent dans deux espaces différents, donc la diagonalisation de $\hat{H}_{\text{tot}}$ peut se faire dans une base propre commune des opérateurs $\hat{H}_{\text{cm}}$ et $\hat{H}$. 

        En utilisant les états propres $\ket{\vec{P}}$ de l’opérateur impulsion totale, on peut écrire les fonctions propres de l’hamiltonienn total sous la forme 
        \[ \Phi(\vec{R}, \vec{r}) = \frac{e^{i \frac{\vec{P} \cdot \vec{R}}{\hbar}}}{(2\pi \hbar)^{\frac{3}{2}}} \psi(\vec{r}) \]   
        où $\psi(\vec{r})$ est une fonction propre de $\hat{H}$ associée à la valeur propre $E$. L’énergie totale sera alors égale à $\frac{P^2}{2M} + E$. 
    \end{demo}

    On peut donc, pour l’étude du système, écrire l’hamiltonien du système sous la forme
    \[ \hat{H}_{\text{tot}} = \frac{\hat{P}^2}{2M} + \frac{\hat{p}^2}{2 \mu} + V(\hat{r}) \]   
    où $\hat{\vec{P}} = \hat{\vec{p}}_1 + \hat{\vec{p}}_2$ est l’impulsion totale du système et  $\hat{\vec{p}} = \frac{m_1 \hat{\vec{p}}_2 - m_2 \hat{\vec{p_1}}}{M}$ est l’impulsion de la particule fictive. Le mouvement du centre de masse étant plus aisé à étudier, on s’intéresse surtout au mouvement de la particule fictive dans $\mathcal{L}^2(\mathbf{R}^3)$ associé à l’hamiltonien 
    \[ \hat{H} = \frac{\hat{p}^2}{2\mu} + V(\hat{r}) \] 
    Le problème à deux corps se ramène donc de façon général à celui d’un seul corps dans un potentiel central. Dans le cas de l’atome d’hydrogène, il suffit donc de remplacer la masse de l’électron $m_e$ par la masse réduite $\mu$. Compte tenu des ordres de grandeurs ($\frac{m_p}{m_e} \approx 2000$), on peut considérer en première approximation directement la masse de l’électron, et supposer que le proton est fixe à l’origine.

    Cherchons donc les fonctions propres $\psi(r, \theta, \varphi)$ de l’hamiltonien $\hat{H}$ de la particule fictive. D’après l’équation de Schödinger indépendante du temps $\hat{H} \ket{\psi} = E \ket{\psi}$, 
    \[ \left(-\frac{\hbar^2}{2 \mu} \Delta + V(r)\right) \psi(r, \theta, \varphi) = E \psi(r, \theta, \varphi) \]   
    où le laplacien s’écrit en coordonnées sphériques selon la relation 
    \begin{align*}
        \Delta 
        &= \frac{1}{r} \frac{\partial^2}{\partial r^2} r + \frac{1}{r^2} \left(\frac{1}{\sin(\theta)} \frac{\partial }{\partial \theta} \left(\sin(\theta) \frac{\partial }{\partial \theta} \right) + \frac{1}{\sin^2(\theta)} \frac{\partial^2 }{\partial \varphi^2} \right) \\
        &= \frac{1}{r} \frac{\partial^2}{\partial r^2} r - \frac{\hat{L}^2}{\hbar^2 r^2} \\
    \end{align*} 
    où $\hat{L}^2$ est l’opérateur carré associé au moment cinétique orbital de la particule. On peut finalement réécrire l’équation de Schödinger comme 
    \[ \left(\underset{\textcolor{couleurmatiere}{E_{c, \text{trans}}}}{\frac{-\hbar^2}{2\mu r} \frac{\partial^2}{\partial r^2} r + V(r)} - \underset{\textcolor{couleurmatiere}{E_{c, \text{rot}}}}{\frac{\hat{L}^2}{2 \mu r^2}}\right) \psi(r, \theta, \varphi) = E \psi(r, \theta, \varphi) \]   
    On peut vérifier explicitement que l’hamiltonien commute avec le moment cinétique, comme prévu par l’invariance par rotation. 

    Cherchons maintenant une base propre commune aux trois observables $\hat{H}$, $\hat{L}^2$ et $\hat{L}_z$. Nous avons déjà vu que cette base propre correspond à une partie radiale et une harmoniques sphérique, de la forme
    \[ \psi(r, \theta, \varphi) = R(r) Y_{\ell,m}(\theta, \varphi) \]   
    où la partie radiale reste à déterminer. On sait déjà que 
    \[ \hat{L}^2 \psi(r, \theta, \varphi) = \ell(\ell + 1) \hbar^2 \psi(r, \theta, \varphi) \]   
    Sachant qu’il n’y a plus d’opérateur portant sur la partie angulaire dans l’équation de Schrödinger, on peut simplifier l’équation par $Y_{\ell,m}(\theta, \varphi)$, et en multipliant par $r$, on obtient
    \[ \left(-\frac{\hbar^2}{2\mu}\frac{\text{d}^2}{\text{d}r^2} + \frac{\ell(\ell+1)\hbar^2}{2\mu r^2} + V(r)\right) u(r) = E u(r) \]  
    où $u(r) = r R(r)$. Finalement, on a réussi à se ramener à une équation de Schrödinger indépendante du temps à une dimension 
    \[ \left(-\frac{\hbar^2}{2\mu} \frac{\text{d}^2}{\text{d}r^2} + V_{\text{eff}, \ell}(r)\right)u(r) = Eu(r) \]   
    correspondant au mouvement d’une particule de masse $\mu$ dans le potentiel 
    \[ V_{\text{eff}, \ell}(r) = \underset{\textcolor{couleurmatiere}{\text{potentiel central}}}{V(r)} + \underset{\textcolor{couleurmatiere}{E_{c,\text{rot}}}}{\frac{\ell(\ell + 1) \hbar^2}{2 \mu r^2}} \]  
    La fonction $u(r)$ sur laquelle porte l’équation est telle que $\abs{u(r)}^2 dr$ est la probabilité que la particule se trouve entre deux sphères concentriques de rayons $r$ et $r+ dr$.

        \subsubsection{Cas du potentiel coulombien}

        Appliquons la démarche exposée dans un cadre général à l’atome d’hydrogène, cas correspondant au potentiel coulombien $V(r) = \frac{-e^2}{r}$. On doit donc résoudre l’équation aux valeurs propres $\hat{H}_{\ell} \ket{u} = E \ket{u}$, où 
        \[ \hat{H}_{\ell} = - \frac{\hbar^2}{2\mu} \frac{\text{d}^2}{\text{d}r^2} - \frac{e^2}{r} + \frac{\ell(\ell + 1)\hbar^2}{2 \mu r^2} \]   
        où $\mu \approx m_e$. 

        \begin{omed}{États liés de l’atome d’hydrogène}
            \begin{itemize}
                \item Les niveaux d’énergie sont repérés par un entier positif $n$ et ont pour valeur
                \[ E_n = - \frac{E_I}{n^2} \]   
                où $E_I = \frac{e^2}{2 a_1} = \frac{\mu e^4}{2 \hbar^2} \approx 13,6\text{ev}$ est l’énergie d’ionisation de l’hydrogène et où $a_1 = \frac{\hbar^2}{\mu e^2} \approx 0,053\text{nm}$ est le rayon de Bohr.
                \item Le niveau $E_n$ est dégénéré $n^2$ fois.
                \item Les états propres $\ket{n,\ell,m}$ sont associés aux fonctions d’onde
                \[ \psi_{n,\ell,m}(r, \theta, \varphi) = R_{n, \ell}(r)Y_{\ell, m}(\theta, \varphi) \]   
                où $Y_{\ell, m}(\theta, \varphi)$ est une harmonique sphérique et où $R_{n, \ell}(r) = \frac{u_{n, \ell}(r)}{r}$ correspond au produit d’un polynôme de degré $n'$ par $r^{\ell} \exp\left(-\frac{r}{n a_1}\right)$ l’entier $n' = n - \ell - 1$ est appelé le nombre quantique radial.
                \item Les zéros du polynôme de degré $n'$ étant tous strictement positifs, la fonction $R_{n, \ell}(r)$ s’annule exactement $n'$ fois dans l’intervalle $\intOO{0}{\infty}$, conformément au théorème de Sturm-Liouville.
                \item La fonction d’onde de l’état fondamental s’écrit 
                \[ \psi_{1,0,0}(\vec{r}) = \frac{e^{-\frac{r}{a_1}}}{\sqrt{\pi a_1^3}} \]
            \end{itemize}
        \end{omed}

        \begin{demo}{Résolution}
            On utilisera comme unité de longueur le rayon de Bohr $a_1$ et comme unité d’énergie $E_I$. On introduit donc les variables sans dimension $\rho = \frac{r}{a_1}$ et $\varepsilon = \frac{E}{E_I}$, ce qui nous donne l’équation de Schrödinger 
            \[ \hat{h}_{\ell} u(\rho) = \left(-\frac{\text{d}^2}{\text{d}\rho^2}  + \frac{\ell(\ell + 1)}{\rho^2} - \frac{2}{\rho}\right)u(\rho) = \varepsilon u(\rho) \]

            On introduit l’opérateur différentiel $\hat{a}_{\ell}$ et son adjoint $\hat{a}_{\ell}^{\dagger}$ 
            \begin{align*}
                \hat{a}_{\ell} &= \frac{1}{\ell + 1} - \frac{\ell + 1}{\rho} - \frac{\text{d}}{\text{d}\rho} \\
                \hat{a}_{\ell}^{\dagger} &= \frac{1}{\ell + 1} - \frac{\ell + 1}{\rho} + \frac{\text{d}}{\text{d}\rho} \\
            \end{align*}
            L’effet de l’opération $\dagger$ sur $\hat{a}_{\ell}$ peut se comprendre par le fait que $\hat{p}$ est une observable donc auto-adjoint, ce qui nous donne que $- i \frac{\text{d}}{\text{d}\rho}$ est auto-adjoint, d’où $\frac{\text{d}}{\text{d}\rho}^{\dagger} = - \frac{\text{d}}{\text{d}\rho}$. Après calculs, on remarque que 
            \begin{align*}
                \hat{a}_{\ell} \hat{a}_{\ell}^{\dagger} &= \hat{h}_{\ell} + \frac{1}{(\ell + 1)^2} \\
                \hat{a}_{\ell}^{\dagger} \hat{a}_{\ell} &= \hat{h}_{\ell + 1} + \frac{1}{(\ell + 1)^2} \\  
            \end{align*}
            On peut ainsi réécrire l’équation aux états propres par 
            \[ \left(\hat{a}_{\ell} \hat{a}_{\ell}^{\dagger} - \frac{1}{(\ell + 1)^2}\right) \ket{u} = \varepsilon \ket{u} \]   
            On peut faire agir l’opérateur $\hat{a}_{\ell}^{\dagger}$ sur cette équation, ce qui donne 
            \begin{align*}
                \left( \hat{a}_{\ell}^{\dagger} \hat{a}_{\ell} - \frac{1}{(\ell + 1)^2}\right)\hat{a}_{\ell}^{\dagger} \ket{u} = \varepsilon \hat{a}_{\ell}^{\dagger} \ket{u}
            \end{align*}
            soit $\hat{h}_{\ell + 1} \hat{a}_{\ell}^{\dagger} \ket{u} = \varepsilon \hat{a}_{\ell}^{\dagger} \ket{u}$. Le vecteur $\hat{a}_{\ell}^{\dagger} \ket{u}$ est donc soit nul soit vecteur propre de $\hat{h}_{\ell + 1}$ pour la même valeur propre $\varepsilon$. De même, en faisant agir l’opérateur $\hat{a}_{\ell - 1}$ sur cette équation, on obtient 
            \[ \left(\hat{a}_{\ell - 1} \hat{a}_{\ell - 1}^{\dagger} - \frac{1}{\ell^2} \right) \hat{a}_{\ell - 1} \ket{u} = \varepsilon \hat{a}_{\ell - 1} \ket{u} \]
            soit $\hat{h}_{\ell - 1} \hat{a}_{\ell - 1} \ket{u} \varepsilon \hat{a}_{\ell - 1} \ket{u}$. Le vecteur $\hat{a}_{\ell - 1} \ket{u}$ est donc soit nul soit vecteur propre de $\hat{h}_{\ell - 1}$ pour la même valeur propre $\varepsilon$. Par ailleurs, le carré de la norme de $\hat{a}_{\ell}^{\dagger} \ket{u}$ s’écrit 
            \begin{align*}
                \nor{\hat{a}_{\ell}^{\dagger} \ket{u}}^2
                &= \bra{u} \hat{a}_{\ell} \hat{a}_{\ell}^{\dagger} \ket{u} \\
                &= \bra{u} \left(\hat{h}_{\ell} + \frac{1}{(\ell + 1)^2}\right) \ket{u} \\
                &= \varepsilon + \frac{1}{(\ell + 1)^2}
                &\geq 0
            \end{align*}
            On a donc la condition 
            \[ \varepsilon \geq - \frac{1}{(\ell + 1)^2} \]

            Employons désormais un raisonnement similaire à celui utilisé par Dirac pour l’oscillateur harmonique ou par Cartan pour le moment cinétique : pour une valeur donnée de $\varepsilon < 0$, on peut appliquer successivement $\hat{a}_{\ell}^{\dagger}$, $\hat{a}_{\ell + 1}^{\dagger}$, etc. jusqu’à ce que l’inégalité du dessus ne soit plus vérifiée. Ceci étant physiquement impossible, il existe $\ell_{\text{max}}$ tel que $\ket{u}$ soit vecteur propre de $\hat{h}_{\ell_{\text{max}}}$ pour la valeur propre $\varepsilon$ mais que $\hat{a}_{\ell_{\text{max}}}^{\dagger} \ket{u}$ soit le vecteur nul. En considérant la norme de ce vecteur et en posant $n = \ell_{\text{max}} + 1$, on obtient donc que 
            \[ \varepsilon = \varepsilon_n = -\frac{1}{n^2} \]   
            Dans le cas où $\ell = \ell_{\text{max}} = n-1$, on peut écrire $\hat{a}_{n-1}^{\dagger}\ket{u} = 0$ d’où 
            \begin{align*}
                \left(\frac{1}{n} - \frac{n}{\rho} + \frac{\text{d}}{\text{d}\rho} \right) u(\rho) &= 0 \\
                \textit{i.e.} \quad \frac{\text{d}u}{u} = - \left(\frac{1}{n} + \frac{n}{\rho}\right) \text{d}\rho
            \end{align*}
            Cette équation différentielle, une fois intégrée, donne que 
            \[ u_{n, n-1}(\rho) \propto \rho^n \exp\left(-\frac{\rho}{n}\right) \]
            Pour trouver les autres fonctions radiales, on utilise la relation 
            \[ \ket{u_{n,\ell - 1}} = \frac{\hat{a}_{\ell - 1} \ket{u_{n, \ell}}}{\nor{{\hat{a}_{\ell - 1} \ket{u_{n, \ell}}}}} \]
            le dénominateur étant non nul car $\ell < n$ et $\nor{{\hat{a}_{\ell - 1} \ket{u_{n, \ell}}}}^2 = \frac{1}{\ell^2} - \frac{1}{n^2}$.

            Le résultat montré sur le degré du polynôme se montre par récurrence, sachant que $u_{n, n-1}(\rho)$ est constitué d’un polynôme de degré $n' = n - \ell - 1 = 0$ et que l’action de $\hat{a}_{\ell - 1} = \left(\frac{1}{\ell} - \frac{\ell}{\rho} - \frac{\text{d}}{\text{d}\rho} \right)$ sur $u_{n,\ell}(\rho) = P(\rho) \rho^{\ell + 1} \exp\left(-\frac{\rho}{n}\right)$ donne un polynôme de degré supérieur $n' + 1 = n - (\ell - 1) - 1$.
        \end{demo}

        Un des résultats remarquables est l’apparition de l’entier $n$ strictement positif, appelé \textbf{nombre quantique principal}, qui détermine entièrement la valeur de l’énergie $E_n$. À énergie constante, on peut trouver différentes valeurs diverses du moment cinétique puisque $\ell \in \intEN{0}{n-1}$. La dégénérescence du niveau $E_n$ est donc égale à 
        \[ g_n = \sum_{\ell = 0}^{n-1} \left(2\ell + 1\right) = n^2 \]   
        La projection du moment cinétique sur l’axe $z$ prendre enfin la valeur $m \hbar$ où $m \in \intEN{-\ell}{\ell}$, ce qui montre que le moment cinétique ne pourra jamais atteindre la valeur $n\hbar$ utilisée pour établir le résultat du modèle de Bohr. 

    \subsection{Interaction spin-orbite}

        Lors de l’étude précédente des états liés de l’atome d’hydrogène, nous n’avons pas tenu compte du degré de liberté de spin de l’électron. La conséquence de son introduction est \textit{a priori} de multiplier par deux la dégénérescence de tous les niveaux d’énergie identifiés, les portant à $2 n^2$ pour le niveau d’énergie $E_n$. On aurait ainsi 
        \[ \hat{H}_0 \ket{n,\ell,m} \otimes \ket{\pm} = E_n \ket{n, \ell, m} \otimes \ket{\pm} \] 

        Toutefois, l’hamiltonien $\hat{H}_0$ ne décrit pas parfaitement l’état du système complet et il faut y ajouter un terme supplémentaire faisant intervenir explicitement le spin de l’électron en raison d’effets relativistes. On admet que le terme correctif, dit d’\textbf{interaction spin-orbite}, à apporter est :
        \[ \hat{W}_{\text{SO}} = \frac{1}{2 m_e^2 c^2} \frac{1}{\hat{r}} \frac{\text{d}V}{\text{d}r}(\hat{r}) \hat{\vec{L}} \cdot \hat{\vec{S}}_e \]

        En ordre de grandeurs, on remarque que 
        \[ \frac{W_{\text{SO}}}{E_I} \approx \frac{e^4}{\hbar^2 c^2} = \alpha^2 \approx 10^{-4} \]
        où $\alpha$ est la constante de structure fine, ce qui signifie que l’on peut traiter cette variation par la méthode des perturbations. Il est notable que la présence de ce terme $\hat{W}_{\text{SO}}$ implique que le moment cinétique orbital $\hat{\vec{L}}$ ne commute plus avec $\hat{H}$. Il faut maintenant considérer le moment cinétique total $\hat{\vec{J}} = \hat{\vec{L}} + \hat{\vec{S}}_e$ qui correspond au véritable générateur infinitésimal des rotations du système complet. On peut vérifier que 
        \[ \hat{\vec{L}} \cdot \hat{\vec{S}}_e = \frac{1}{2} \left(\hat{J}^2 - \hat{L}^2 - \hat{S}^2_e\right) \]   
        Cette expression montre explicitement que $\hat{W}_{\text{SO}}$ commute avec $\hat{\vec{J}}$, è. Compte tenu des résultats sur l’addition des moments cinétiques, on sait que les valeurs propres de $\hat{J}^2$ seront de la forme $j(j+1)\hbar^2$ avec $j = \ell \pm \frac{1}{2}$.

        Dans le cas de l’atome d’hydrogène, le niveau fondamental $1s$ de l’hamiltonien $\hat{H}_0$ est doublement dégénéré, et une base de celui c-ci peut s’écrire $\left\{\ket{1,0,0} \otimes \ket{\pm}\right\}$. Ces deux états correspondant à $\ell = 0$ sont vecteurs propres de $\hat{L}_x$, $\hat{L}_y$ et $\hat{L}_z$ pour la valeur propre $0$, donc $\hat{W}_0 \ket{1,0,0} \otimes \ket{\pm} = 0$. L’interaction spin-orbite est sans effet sur ces états. Il faudrait de plus considérer le couplage avec le spin du noyau pour observer une levée de dégénérescence, phénomène que l’on appelle \textbf{structure hyperfine}.

    \subsection{Structure fine de l’atome d’hydrogène}

        Dans le cas de l’atome d’hydrogène, le noyau est un proton, de spin 1/2. Ce degré de liberté supplémentaire, une fois considéré, fait apparaître une interaction entre les dipôles magnétiques du proton et de l’électron :
        \[ W_{HF} = \frac{\mu_0}{3\pi r^3} \left(\vec{\mu}_e \cdot \vec{\mu}_p - 3 (\vec{\mu}_e \cdot \vec{u})(\vec{\mu}_p \cdot \vec{u})\right) - \frac{2\mu_0}{3} \vec{\mu}_e \cdot \vec{\mu}_p \delta(\vec{r}) \]
        où $\mu_0$ est la perméabilité magnétique du vide, tandis que $\vec{\mu}_e$ et $\vec{\mu}_p$ sont le moments magnétiques de l’électron et du proton. $\vec{u}$ est un vecteur unitaire porté par le vecteur $\vec{r}$ séparant les deux dipôles. Le \textbf{terme de contact} proportionnelle à la fonction de Dirac $\delta(\vec{r})$ est essentiel car la probabilité que l’électron et le proton soient au même point est non nulle. En ordres de grandeur, $\frac{W_{HF}}{E_I} \approx 10^{-6}$ donc on peut tout à fait traiter ce problème par méthode des perturbations.

        Considérons le cas du niveau fondamental de l’atome d’hydrogène, d’énergie $E_1 = - E_I$, et correspondant aux états de la forme $\ket{1,0,0} \otimes \ket{e : \pm} \otimes \ket{p : \pm}$. 
        
        \begin{omed}{Opérateur $\hat{H}_1$ dans l’état fondamental}
            La restriction de l’opérateur $\hat{W}_{\text{HF}}$ sur l’espace des spins $\left\{\ket{e : \pm, p : \pm}\right\}$ est 
            \[ \hat{H}_1 = A \left(\frac{1}{2} \frac{\hat{S}^2}{\hbar^2} - \frac{3}{4}\hat{I}\right) \] 
            où $A \approx 5.87 \, \mu\text{eV}$. 
        \end{omed}

        \begin{demo}{Preuve}
            Rappelons l’expression de $\hat{W}_{\text{HF}}$, que l’on notera ici 
            \[ \hat{W} := \frac{\mu_0}{4\pi} \left(\frac{\hat{\vec{\mu}}_e \cdot \hat{\vec{\mu}}_p}{\hat{r}^3} - 3 \frac{(\hat{\vec{\mu}}_e \cdot \hat{\vec{r}}) (\hat{\vec{\mu}}_p \cdot \hat{\vec{r}})}{\hat{r}^5}\right) - \frac{2\mu_0}{3} \hat{\vec{\mu}}_e \cdot \hat{\vec{\mu}}_p \delta(\vec{r}) \]
            Les opérateurs $\hat{\vec{\mu}}_{e,p}$ agissent sur les spins, alors que l’opérateur $\hat{\vec{r}}$ agit sur la partie spaciale. On définit l’opérateur $\hat{H_1}$, qui agit sur l’espace des spins, par ses éléments matriciels : 
            \[ \bra{1,0,0,\sigma_e, \sigma_p} \hat{W} \ket{1,0,0, \sigma_e', \sigma_p'} = \bra{\sigma_e, \sigma_p} \hat{H}_1 \ket{\sigma_e', \sigma_p'} \]    
            Autrement dit, pour définir $\hat{H}_1$, on prend $\hat{W}$, puis on « intègre » sur sa partie spaciale en moyennant sur les états orbitaux $\ket{1,0,0}$ :
            \[ \hat{H}_1 = \frac{\mu_0}{4\pi} \left(\left<\frac{1}{r^3}\right>\hat{\vec{\mu}}_e \cdot \hat{\vec{\mu}}_p - 3\left<\frac{(\hat{\vec{\mu}}_e \cdot \vec{r}) (\hat{\vec{\mu}}_p \cdot {\vec{r}})}{r^5}\right> \right) - \frac{2\mu_0}{3} \hat{\vec{\mu}}_e \cdot \hat{\vec{\mu}}_p \left<\delta(\vec{r})\right> \]
            Calculons les moyennes intervenant dans cette formule. Par symétrie, on sait que les produits associés à deux composantes cartésiennes différentes (comme $\left< xy \right>$) sont nuls. De même, par symétrie, $\left<x^2 \right> = \left<y^2 \right> = \left<z^2 \right> = \frac{1}{3} \left<r^2 \right>$. Ainsi,
            \begin{align*}
                \left<\frac{(\hat{\vec{\mu}}_e \cdot \vec{r}) (\hat{\vec{\mu}}_p \cdot {\vec{r}})}{r^5}\right> 
                &= \hat{\mu}_{ex} \hat{\mu}_{px}\left<\frac{x^2}{r^5}\right> + \hat{\mu}_{ey} \hat{\mu}_{py} \left<\frac{y^2}{r^5}\right>+ \hat{\mu}_{ez} \hat{\mu}_{pz} \left<\frac{z^2}{r^5}\right> \\
                &= \frac{1}{3} \left<\frac{1}{r^3}\right> \hat{\vec{\mu}}_e \cdot \hat{\vec{\mu}}_p \\
            \end{align*}
            Cette moyenne s’annule donc avec l’autre membre de la parenthèse, et il ne reste qu’à calculer 
            \[ \left<\delta(\vec{r})\right> = \iiint \abs{\psi_{1s}(\vec{r})}^2\delta(\vec{r}) d^3r = \abs{\psi_{1s}(0)}^2 \]
            Sachant que l’expression de la partie radiale de la fonction d’onde du niveau fondamental s’écrit $\psi_{1s}(\vec{r}) = \sqrt{\frac{1}{\pi a_1^3}} \exp\left(-\frac{r}{a_1}\right)$, on obtient finalement 
            \[ \hat{H_1} = - \frac{2\mu_0}{3\pi a_1^3} \hat{\vec{\mu}}_e \cdot \hat{\vec{\mu}}_p \]
            On conclut en utilisant que $\hat{\vec{\mu}}_{e,p} = \gamma_{e,p} \hat{\vec{S}}_{e,p}$.

            Or on a que $\hat{S}^2 = \hat{S}^2_e + \hat{S}^2_p + 2 \hat{\vec{S}}_e \cdot \hat{\vec{S}}_p$, ce qui nous permet d’écrire que 
            \[ \hat{\vec{S}}_e \cdot \hat{\vec{S}}_p = \frac{1}{2} \left(\hat{S}^2 - \hat{S}^2_e - \hat{S}^2_p\right) \]
            Les particules $e$ et $p$ étant de spin 1/2, $\hat{S}_e^2 = \hat{S}_p^2 = \frac{3}{4} \hbar^2 \hat{I}$ et donc 
            \[ \hat{H}_1 = A \left(\frac{1}{2} \frac{\hat{S}^2}{\hbar^2} - \frac{3}{4}\hat{I}\right) \]  
        \end{demo}
          
        L’hamiltonien de structure hyperfine s’exprime donc directement en fonction de l’observable $\hat{S}^2$. Le cas de l’addition de deux spins 1/2 ayant déjà été traité dans ce poly, la base couplée constituant une base propre de l’hamiltonien, on peut écrire directement que 
        \[ \hat{H}_1 \ket{s,m} = A \left(\frac{s(s+1)}{2} - \frac{3}{4}\right) \ket{s,m} \]   
        On a donc $\hat{H}_1 \ket{1,m} = \frac{A}{4} \ket{1,m}$ et $\hat{H}_1 \ket{0,0} = - \frac{3A}{4} \ket{0,0}$. Le niveau fondamental de l’atome d’hydrogène est ainsi clivé en deux sous-niveaux. Un sous-niveau singulet, non dégénéré, qui devient le véritable état fondamental du système, ainsi qu’un sous-niveau triplet, dégénéré trois fois. L’amplitude du clivage est de l’ordre de $A \approx 5.87 \, \mu\text{eV}$, ce qui correspond à une fréquence de $\nu \approx 1.42 \, \text{GHz}$ et une longueur d’onde de $\lambda \approx 21 \, \text{cm}$




    

    

    